% Options for packages loaded elsewhere
\PassOptionsToPackage{unicode}{hyperref}
\PassOptionsToPackage{hyphens}{url}
\PassOptionsToPackage{dvipsnames,svgnames,x11names}{xcolor}
\documentclass[
  11pt,
  a4paper,
]{article}
\usepackage{xcolor}
\usepackage[margin=1in]{geometry}
\usepackage{amsmath,amssymb}
\setcounter{secnumdepth}{-\maxdimen} % remove section numbering
\usepackage{iftex}
\ifPDFTeX
  \usepackage[T1]{fontenc}
  \usepackage[utf8]{inputenc}
  \usepackage{textcomp} % provide euro and other symbols
\else % if luatex or xetex
  \usepackage{unicode-math} % this also loads fontspec
  \defaultfontfeatures{Scale=MatchLowercase}
  \defaultfontfeatures[\rmfamily]{Ligatures=TeX,Scale=1}
\fi
\usepackage{lmodern}
\ifPDFTeX\else
  % xetex/luatex font selection
  \ifLuaTeX
    \usepackage{luaotfload}
    \directlua{luaotfload.add_fallback("mainfontfallback",{
      "NotoSerifBengali:mode=harf"
    })}
  \fi
  \setmainfont[,RawFeature={fallback=mainfontfallback}]{Latin Modern Roman}
  \ifLuaTeX
    \usepackage{luaotfload}
    \directlua{luaotfload.add_fallback("sansfontfallback",{
      "NotoSansBengali:mode=harf"
    })}
  \fi
  \setsansfont[,RawFeature={fallback=sansfontfallback}]{Latin Modern Sans}
  \ifLuaTeX
    \usepackage{luaotfload}
    \directlua{luaotfload.add_fallback("monofontfallback",{
      "NotoSansBengali:mode=harf"
    })}
  \fi
  \setmonofont[,RawFeature={fallback=monofontfallback}]{DejaVu Sans Mono}
\fi
% Use upquote if available, for straight quotes in verbatim environments
\IfFileExists{upquote.sty}{\usepackage{upquote}}{}
\IfFileExists{microtype.sty}{% use microtype if available
  \usepackage[]{microtype}
  \UseMicrotypeSet[protrusion]{basicmath} % disable protrusion for tt fonts
}{}
\makeatletter
\@ifundefined{KOMAClassName}{% if non-KOMA class
  \IfFileExists{parskip.sty}{%
    \usepackage{parskip}
  }{% else
    \setlength{\parindent}{0pt}
    \setlength{\parskip}{6pt plus 2pt minus 1pt}}
}{% if KOMA class
  \KOMAoptions{parskip=half}}
\makeatother
\usepackage{color}
\usepackage{fancyvrb}
\newcommand{\VerbBar}{|}
\newcommand{\VERB}{\Verb[commandchars=\\\{\}]}
\DefineVerbatimEnvironment{Highlighting}{Verbatim}{commandchars=\\\{\}}
% Add ',fontsize=\small' for more characters per line
\usepackage{framed}
\definecolor{shadecolor}{RGB}{248,248,248}
\newenvironment{Shaded}{\begin{snugshade}}{\end{snugshade}}
\newcommand{\AlertTok}[1]{\textcolor[rgb]{0.94,0.16,0.16}{#1}}
\newcommand{\AnnotationTok}[1]{\textcolor[rgb]{0.56,0.35,0.01}{\textbf{\textit{#1}}}}
\newcommand{\AttributeTok}[1]{\textcolor[rgb]{0.13,0.29,0.53}{#1}}
\newcommand{\BaseNTok}[1]{\textcolor[rgb]{0.00,0.00,0.81}{#1}}
\newcommand{\BuiltInTok}[1]{#1}
\newcommand{\CharTok}[1]{\textcolor[rgb]{0.31,0.60,0.02}{#1}}
\newcommand{\CommentTok}[1]{\textcolor[rgb]{0.56,0.35,0.01}{\textit{#1}}}
\newcommand{\CommentVarTok}[1]{\textcolor[rgb]{0.56,0.35,0.01}{\textbf{\textit{#1}}}}
\newcommand{\ConstantTok}[1]{\textcolor[rgb]{0.56,0.35,0.01}{#1}}
\newcommand{\ControlFlowTok}[1]{\textcolor[rgb]{0.13,0.29,0.53}{\textbf{#1}}}
\newcommand{\DataTypeTok}[1]{\textcolor[rgb]{0.13,0.29,0.53}{#1}}
\newcommand{\DecValTok}[1]{\textcolor[rgb]{0.00,0.00,0.81}{#1}}
\newcommand{\DocumentationTok}[1]{\textcolor[rgb]{0.56,0.35,0.01}{\textbf{\textit{#1}}}}
\newcommand{\ErrorTok}[1]{\textcolor[rgb]{0.64,0.00,0.00}{\textbf{#1}}}
\newcommand{\ExtensionTok}[1]{#1}
\newcommand{\FloatTok}[1]{\textcolor[rgb]{0.00,0.00,0.81}{#1}}
\newcommand{\FunctionTok}[1]{\textcolor[rgb]{0.13,0.29,0.53}{\textbf{#1}}}
\newcommand{\ImportTok}[1]{#1}
\newcommand{\InformationTok}[1]{\textcolor[rgb]{0.56,0.35,0.01}{\textbf{\textit{#1}}}}
\newcommand{\KeywordTok}[1]{\textcolor[rgb]{0.13,0.29,0.53}{\textbf{#1}}}
\newcommand{\NormalTok}[1]{#1}
\newcommand{\OperatorTok}[1]{\textcolor[rgb]{0.81,0.36,0.00}{\textbf{#1}}}
\newcommand{\OtherTok}[1]{\textcolor[rgb]{0.56,0.35,0.01}{#1}}
\newcommand{\PreprocessorTok}[1]{\textcolor[rgb]{0.56,0.35,0.01}{\textit{#1}}}
\newcommand{\RegionMarkerTok}[1]{#1}
\newcommand{\SpecialCharTok}[1]{\textcolor[rgb]{0.81,0.36,0.00}{\textbf{#1}}}
\newcommand{\SpecialStringTok}[1]{\textcolor[rgb]{0.31,0.60,0.02}{#1}}
\newcommand{\StringTok}[1]{\textcolor[rgb]{0.31,0.60,0.02}{#1}}
\newcommand{\VariableTok}[1]{\textcolor[rgb]{0.00,0.00,0.00}{#1}}
\newcommand{\VerbatimStringTok}[1]{\textcolor[rgb]{0.31,0.60,0.02}{#1}}
\newcommand{\WarningTok}[1]{\textcolor[rgb]{0.56,0.35,0.01}{\textbf{\textit{#1}}}}
\usepackage{longtable,booktabs,array}
\usepackage{calc} % for calculating minipage widths
% Correct order of tables after \paragraph or \subparagraph
\usepackage{etoolbox}
\makeatletter
\patchcmd\longtable{\par}{\if@noskipsec\mbox{}\fi\par}{}{}
\makeatother
% Allow footnotes in longtable head/foot
\IfFileExists{footnotehyper.sty}{\usepackage{footnotehyper}}{\usepackage{footnote}}
\makesavenoteenv{longtable}
\setlength{\emergencystretch}{3em} % prevent overfull lines
\providecommand{\tightlist}{%
  \setlength{\itemsep}{0pt}\setlength{\parskip}{0pt}}
% paper_preamble.tex — shared LaTeX preamble for the PCMA/SI paper set.
%
% Compiled with lualatex (NOT pdflatex) because several papers contain
% Bengali script: the CLL paper realises one cognitive state into
% English, Bengali and Banglish, so the Bengali output is EVIDENCE and
% cannot be transliterated away without destroying the point being made.
%
% Font selection is NOT done here. pandoc >= 3 has native
% `mainfontfallback` / `monofontfallback` variables that emit the
% luaotfload fallback registration in the correct position relative to
% unicode-math; doing it by hand in this header lands after unicode-math
% and breaks font loading outright. See tools/build_papers.sh.

\usepackage{amsmath}
\usepackage{amssymb}
\usepackage{mathtools}   % \xleftrightarrow etc. used in cll.md §5
\usepackage{booktabs}

% ---------------------------------------------------------------------
% Unicode symbols appearing in PROSE rather than inside $...$.
%
% The papers were written as markdown and read as markdown, so they use
% real characters — "the ω-boost", "α-child", "ΔR coherence" — instead of
% LaTeX macros. Under unicode-math those codepoints are math-mode-only
% and abort the build ("\symbballowed only in math mode"). Mapping each
% to \ensuremath keeps the markdown readable on GitHub AND compilable,
% which is the whole point of keeping one source.
% ---------------------------------------------------------------------
\usepackage{newunicodechar}

% Greek, lower case
\newunicodechar{α}{\ensuremath{\alpha}}
\newunicodechar{β}{\ensuremath{\beta}}
\newunicodechar{γ}{\ensuremath{\gamma}}
\newunicodechar{δ}{\ensuremath{\delta}}
\newunicodechar{ε}{\ensuremath{\varepsilon}}
\newunicodechar{η}{\ensuremath{\eta}}
\newunicodechar{θ}{\ensuremath{\theta}}
\newunicodechar{κ}{\ensuremath{\kappa}}
\newunicodechar{λ}{\ensuremath{\lambda}}
\newunicodechar{μ}{\ensuremath{\mu}}
\newunicodechar{π}{\ensuremath{\pi}}
\newunicodechar{ρ}{\ensuremath{\rho}}
\newunicodechar{σ}{\ensuremath{\sigma}}
\newunicodechar{τ}{\ensuremath{\tau}}
\newunicodechar{φ}{\ensuremath{\varphi}}
\newunicodechar{χ}{\ensuremath{\chi}}
\newunicodechar{ψ}{\ensuremath{\psi}}
\newunicodechar{ω}{\ensuremath{\omega}}
\newunicodechar{ν}{\ensuremath{\nu}}
\newunicodechar{ξ}{\ensuremath{\xi}}
\newunicodechar{ζ}{\ensuremath{\zeta}}
\newunicodechar{ι}{\ensuremath{\iota}}

% Greek, upper case
\newunicodechar{Δ}{\ensuremath{\Delta}}
\newunicodechar{Φ}{\ensuremath{\Phi}}
\newunicodechar{Σ}{\ensuremath{\Sigma}}
\newunicodechar{Ψ}{\ensuremath{\Psi}}
\newunicodechar{Ω}{\ensuremath{\Omega}}
\newunicodechar{Γ}{\ensuremath{\Gamma}}
\newunicodechar{Λ}{\ensuremath{\Lambda}}
\newunicodechar{Θ}{\ensuremath{\Theta}}
\newunicodechar{Π}{\ensuremath{\Pi}}

% Operators and relations
\newunicodechar{×}{\ensuremath{\times}}
\newunicodechar{·}{\ensuremath{\cdot}}
\newunicodechar{±}{\ensuremath{\pm}}
\newunicodechar{−}{\ensuremath{-}}
\newunicodechar{≈}{\ensuremath{\approx}}
\newunicodechar{≤}{\ensuremath{\leq}}
\newunicodechar{≥}{\ensuremath{\geq}}
\newunicodechar{≠}{\ensuremath{\neq}}
\newunicodechar{∈}{\ensuremath{\in}}
\newunicodechar{∂}{\ensuremath{\partial}}
\newunicodechar{∇}{\ensuremath{\nabla}}
\newunicodechar{√}{\ensuremath{\sqrt{}}}
\newunicodechar{∞}{\ensuremath{\infty}}
\newunicodechar{∝}{\ensuremath{\propto}}
\newunicodechar{∑}{\ensuremath{\sum}}
\newunicodechar{∫}{\ensuremath{\int}}
\newunicodechar{ℝ}{\ensuremath{\mathbb{R}}}
\newunicodechar{ℕ}{\ensuremath{\mathbb{N}}}
\newunicodechar{‖}{\ensuremath{\|}}
\newunicodechar{⊕}{\ensuremath{\oplus}}
\newunicodechar{⊗}{\ensuremath{\otimes}}
\newunicodechar{∀}{\ensuremath{\forall}}
\newunicodechar{∃}{\ensuremath{\exists}}
\newunicodechar{¬}{\ensuremath{\neg}}
\newunicodechar{∧}{\ensuremath{\wedge}}
\newunicodechar{∨}{\ensuremath{\vee}}

% Arrows
\newunicodechar{→}{\ensuremath{\rightarrow}}
\newunicodechar{←}{\ensuremath{\leftarrow}}
\newunicodechar{↔}{\ensuremath{\leftrightarrow}}
\newunicodechar{↑}{\ensuremath{\uparrow}}
\newunicodechar{↓}{\ensuremath{\downarrow}}
\newunicodechar{⇒}{\ensuremath{\Rightarrow}}
\newunicodechar{⇐}{\ensuremath{\Leftarrow}}
\newunicodechar{⇔}{\ensuremath{\Leftrightarrow}}

% Super/subscripts written as literal characters
\newunicodechar{²}{\ensuremath{^{2}}}
\newunicodechar{³}{\ensuremath{^{3}}}
\newunicodechar{⁴}{\ensuremath{^{4}}}
\newunicodechar{⁵}{\ensuremath{^{5}}}
\newunicodechar{⁶}{\ensuremath{^{6}}}
\newunicodechar{₀}{\ensuremath{_{0}}}
\newunicodechar{₁}{\ensuremath{_{1}}}
\newunicodechar{₂}{\ensuremath{_{2}}}
\newunicodechar{₃}{\ensuremath{_{3}}}
\newunicodechar{ᵀ}{\ensuremath{^{\mathsf{T}}}}

% Status glyphs. The papers use emoji as status markers in tables
% (implemented / failed / warning). Mapped to typographic equivalents:
% a colour emoji font in lualatex is unreliable, and the MEANING is what
% carries — a reader needs to see "this one failed", not a red cross.
\newunicodechar{✅}{\ensuremath{\checkmark}}
\newunicodechar{✓}{\ensuremath{\checkmark}}
\newunicodechar{❌}{\ensuremath{\times}}
\newunicodechar{✗}{\ensuremath{\times}}
\newunicodechar{⚠}{\textbf{!}}
\newunicodechar{🧩}{\textbf{[piece]}}
\newunicodechar{📋}{\textbf{[list]}}
\newunicodechar{🟡}{\textbf{[partial]}}
\newunicodechar{🔴}{\textbf{[dead]}}
\newunicodechar{🟢}{\textbf{[live]}}
\newunicodechar{▼}{\ensuremath{\blacktriangledown}}
\newunicodechar{►}{\ensuremath{\blacktriangleright}}
\newunicodechar{□}{\ensuremath{\square}}
\newunicodechar{█}{\rule{0.6em}{0.9ex}}
\newunicodechar{•}{\textbullet}
\newunicodechar{†}{\textdagger}
\newunicodechar{‑}{-}
\newunicodechar{‐}{-}
\newunicodechar{…}{\ldots}
\newunicodechar{।}{\textbar}   % danda: rendered by the Bengali fallback in text
\newunicodechar{⁻}{\ensuremath{^{-}}}

% Zero-width formatting codepoints. These are invisible in an editor but
% are real characters to LaTeX: an emoji written as "⚠️" is TWO
% codepoints (U+26A0 U+FE0F), so declaring only the visible one leaves a
% variation selector with no definition and the build dies on a symbol
% nobody can see. Both are mapped to nothing.
\newunicodechar{^^^^fe0f}{}    % VARIATION SELECTOR-16
\newunicodechar{^^^^200d}{}    % ZERO WIDTH JOINER
\newunicodechar{^^^^200b}{}    % ZERO WIDTH SPACE

\usepackage{longtable}
\usepackage{array}

% Comparison matrices and cost ledgers are wide. Let them shrink to the
% text block rather than overflow the margin.
\usepackage{adjustbox}

% Code blocks: wrap long lines instead of running off the page. Several
% papers quote real C++ with long identifiers.
\usepackage{fvextra}
\fvset{breaklines=true,breakanywhere=true,fontsize=\small}

\setlength{\parskip}{0.5em}
\setlength{\parindent}{0pt}

\usepackage{fancyhdr}
\pagestyle{fancy}
\fancyhf{}
\fancyfoot[C]{\thepage}
\fancyhead[L]{\small\itshape Synthetic Intelligence — research lineage}
\fancyhead[R]{\small\itshape irfan@furylogic.com}
\renewcommand{\headrulewidth}{0.2pt}
\setlength{\headheight}{14pt}
\setlength{\headsep}{18pt}
\usepackage{bookmark}
\IfFileExists{xurl.sty}{\usepackage{xurl}}{} % add URL line breaks if available
\urlstyle{same}
\hypersetup{
  pdftitle={Hierarchical Cognitive Manifold Architecture (HCMA): Sequence Generation via Time-Dependent Energy Minimization},
  pdfauthor={Irfan Mahir --- Independent Researcher, Dhaka, Bangladesh --- irfan@furylogic.com},
  colorlinks=true,
  linkcolor={black},
  filecolor={Maroon},
  citecolor={Blue},
  urlcolor={[HTML]{0B5394}},
  pdfcreator={LaTeX via pandoc}}

\title{Hierarchical Cognitive Manifold Architecture (HCMA): Sequence
Generation via Time-Dependent Energy Minimization}
\author{Irfan Mahir --- Independent Researcher, Dhaka, Bangladesh ---
irfan@furylogic.com}
\date{}

\begin{document}
\maketitle

{
\hypersetup{linkcolor=}
\setcounter{tocdepth}{3}
\tableofcontents
}
\section{Hierarchical Cognitive Manifold Architecture (HCMA): Sequence
Generation via Time-Dependent Energy
Minimization}\label{hierarchical-cognitive-manifold-architecture-hcma-sequence-generation-via-time-dependent-energy-minimization}

\textbf{⚠️ UPDATE: This paper describes the original approach
(hierarchical states + flow field for text generation). The current
implementation evolved into a simpler, more fundamental ``thinking
machine'': a single 256D cognitive state with Hamiltonian reasoning on a
learned potential Φ(c). No text generation, no flow field, no hierarchy.
See the \href{../src/chatbot.cpp}{chatbot} for the current
architecture.}

\textbf{Irfan Mahir}\\
\emph{Independent Researcher, Dhaka, Bangladesh}\\
irfan@furylogic.com

\begin{quote}
\emph{``Thinking and understanding is the art of innovation and
invention.''}
\end{quote}

\begin{center}\rule{0.5\linewidth}{0.5pt}\end{center}

\subsection{Abstract}\label{abstract}

We present the \textbf{Hierarchical Cognitive Manifold Architecture
(HCMA)} , a novel non‑Transformer, non‑autoregressive framework for
sequence generation that overcomes the fundamental limitations of static
potential‐based models. Standard cognitive architectures like P‑CMA
suffer from two critical issues when applied to text generation: (1) the
contrastive Hebbian learning rule collapses hidden units due to a global
scalar teaching signal, and (2) Hamiltonian dynamics in a static
potential settles into a local minimum and cannot traverse a sequence of
states. HCMA addresses both by introducing a \textbf{hierarchical
two‑level cognitive state} (slow and fast) and a \textbf{time‑dependent
driving potential} that propels the system forward through the state
manifold. The slow state encodes high‑level context and modulates the
fast state, while the fast state generates the actual text. The energy
landscape is made time‑dependent via a learned ``flow field'' that
biases the dynamics toward the next text window. Crucially, HCMA retains
all the advantages of P‑CMA---O(1) memory, no tokenization, no
encoders/decoders, and no loss functions---while enabling coherent,
open‑ended sequence generation. We demonstrate the architecture's
effectiveness on synthetic text data, showing that it produces fluent
and diverse outputs without hallucination or catastrophic forgetting.

\begin{center}\rule{0.5\linewidth}{0.5pt}\end{center}

\subsection{1. Introduction}\label{introduction}

Recent attempts to adapt the Predictive Cognitive Manifold Architecture
(P‑CMA) to text generation have uncovered two deep‑seated problems.
First, the contrastive Hebbian update for the potential Φ(c) uses a
global scalar energy signal that is applied uniformly to all hidden
units. This leads to a \textbf{hidden‑layer collapse}: every hidden unit
receives the same gradient direction, preventing specialization. Second,
the Hamiltonian dynamics of a single particle in a static potential
cannot produce a sequence: the particle simply falls into the nearest
local minimum and remains trapped there.

These issues are not mere engineering bugs---they are
\textbf{mathematical consequences} of using a single scalar potential to
both represent ``goodness'' and drive generation. To generate a
sequence, the system must transition from one coherent state to the
next, which requires either a time‑dependent landscape or a mechanism to
inject energy to escape local minima.

We propose a novel architecture that addresses both problems
simultaneously. The core ideas are:

\begin{enumerate}
\def\labelenumi{\arabic{enumi}.}
\item
  \textbf{Hierarchical state decomposition}: Split the cognitive state
  into a slow component (which evolves at the sentence level) and a fast
  component (which evolves at the byte level). The slow state modulates
  the fast state's dynamics, providing a source of temporal structure.
\item
  \textbf{Time‑dependent potential}: Learn a \textbf{flow potential}
  that depends on the current state and an internal ``clock'' or desired
  progress, effectively adding a drifting force that pushes the fast
  state toward the next text window.
\item
  \textbf{Decoupled contrastive Hebbian learning (DCHL)}: Replace the
  global scalar energy signal with \textbf{per-unit contrastive errors}
  \(h^{\text{pos}}_j - h^{\text{neg}}_j\) for the hidden layer and
  independent target projection
  \((0-\Phi^{\text{pos}}) \cdot h^{\text{pos}} + (1-\Phi^{\text{neg}}) \cdot h^{\text{neg}}\)
  for the output layer. Hidden units specialize via self-amplifying
  positive feedback, while the output layer suppresses wrong-polarity
  units via w₂ renormalization.
\end{enumerate}

HCMA remains fully non‑autoregressive, uses no tokens or embeddings, and
requires neither encoders nor decoders. The state itself is the text,
and generation is achieved through energy‑driven evolution with an
explicit forward bias.

\begin{center}\rule{0.5\linewidth}{0.5pt}\end{center}

\subsection{2. The Mathematical
Framework}\label{the-mathematical-framework}

\subsubsection{2.1 Two‑Level State
Representation}\label{twolevel-state-representation}

Let the cognitive state be decomposed into:

\begin{itemize}
\tightlist
\item
  \textbf{Slow state} \(s \in \mathbb{R}^{d_s}\): represents long‑term
  context (e.g., topic, style, narrative arc). It evolves on a coarse
  time scale (e.g., every few sentences).
\item
  \textbf{Fast state} \(f \in \mathbb{R}^{d_f}\): represents the current
  text window (e.g., 256 bytes). It evolves on a fine time scale (e.g.,
  every byte).
\end{itemize}

The full state is \(\mathbf{x} = (s, f)\). The fast state is what we
directly read out as text (via the byte‑to‑float mapping). The slow
state modulates the energy landscape of the fast state.

\subsubsection{2.2 Time‑Dependent
Potential}\label{timedependent-potential}

We define the total energy as:

\[
\Phi(\mathbf{x}, t) = \Phi_f(f, s) + \Phi_s(s) + \Phi_{\text{flow}}(f, s, t)
\]

\begin{itemize}
\tightlist
\item
  \(\Phi_f(f, s)\): energy of the fast state given the slow context. Low
  when the fast window is coherent and consistent with the context.
\item
  \(\Phi_s(s)\): energy of the slow state, ensuring consistency of the
  high‑level structure.
\item
  \(\Phi_{\text{flow}}(f, s, t)\): a \textbf{time‑dependent driving
  potential} that decreases as the system moves along the ``correct''
  sequence. It is designed to push the fast state forward.
\end{itemize}

The flow potential is defined as:

\[
\Phi_{\text{flow}}(f, s, t) = - \alpha \cdot \langle \nabla_f \Phi_f(f, s), \; v_{\text{target}} \rangle \cdot \Delta(t)
\]

where \(v_{\text{target}}\) is a learned vector field pointing toward
the next coherent window, and \(\Delta(t)\) is a monotonically
increasing function of time (e.g., a ramp). This term lowers the energy
in the direction of forward progress, effectively creating a
``gradient'' that drives the state to leave the current minimum.

\subsubsection{2.3 Dynamics}\label{dynamics}

The fast state evolves under Hamiltonian dynamics with an added damping
and a driving term:

\[
\dot{f} = p_f, \quad \dot{p}_f = -\nabla_f \Phi(\mathbf{x}, t) - \gamma p_f + \xi(t)
\]

where \(\gamma\) is a damping coefficient and \(\xi(t)\) is a small
stochastic noise to help escape shallow minima. The slow state evolves
on a much slower timescale:

\[
\dot{s} = p_s, \quad \dot{p}_s = -\nabla_s \Phi(\mathbf{x}, t)
\]

The driving term \(\xi(t)\) and the flow potential ensure that the fast
state does not get stuck.

\subsubsection{2.4 Training via Decoupled Contrastive Hebbian Learning
(DCHL)}\label{training-via-decoupled-contrastive-hebbian-learning-dchl}

Standard contrastive Hebbian learning uses a single scalar teaching
signal \(\Phi\) to update all weights, causing hidden-layer collapse. We
replace this with \textbf{decoupled contrastive Hebbian learning
(DCHL)}: separate update rules for the hidden and output layers, each
using distinct error signals.

\paragraph{Hidden Layer: Per-Unit Contrastive
Hebbian}\label{hidden-layer-per-unit-contrastive-hebbian}

Each hidden unit computes a \textbf{local contrastive error} by
comparing its activation for positive (text) and negative (shuffled)
inputs:

\[
\text{error}_j = h_j^{\text{pos}} - h_j^{\text{neg}}
\]

The input weights of unit \(j\) are updated using only the
\textbf{positive-phase input} \(\mathbf{c}^{\text{pos}}\), creating a
self-amplifying feedback loop:

\[
\Delta W_1[j,:] = \eta_{\text{hidden}} \cdot (h_j^{\text{pos}} - h_j^{\text{neg}}) \cdot \mathbf{c}^{\text{pos}}
\]

\[
\Delta b_{1,j} = \eta_{\text{hidden}} \cdot (h_j^{\text{pos}} - h_j^{\text{neg}})
\]

This rule has three key properties:

\begin{enumerate}
\def\labelenumi{\arabic{enumi}.}
\tightlist
\item
  \textbf{Positive feedback}: Any initial difference between
  \(h^{\text{pos}}\) and \(h^{\text{neg}}\) is amplified. If a unit
  activates more for text, it strengthens its connection to text-like
  patterns; if it activates more for shuffled, it weakens its connection
  to text-like patterns.
\item
  \textbf{Phase-selective coupling}: The update couples to
  \(\mathbf{c}^{\text{pos}}\) only (same-phase input). Using the
  cross-phase input would create homeostatic anti-separation dynamics.
\item
  \textbf{No target information}: The hidden layer never sees the
  targets (0 and 1). All target-driven separation is handled by the
  output layer.
\end{enumerate}

\paragraph{Output Layer: Projection to
Targets}\label{output-layer-projection-to-targets}

The output layer independently projects hidden activations to the target
values:

\[
\Phi = \mathbf{w}_2^T \mathbf{h} + b_2
\]

The output weights are trained to minimize \(\Phi^{\text{pos}} \to 0\)
and \(\Phi^{\text{neg}} \to 1\):

\[
\Delta w_{2,j} = \eta_{\text{output}} \cdot \left( (0 - \Phi^{\text{pos}}) \cdot h_j^{\text{pos}} + (1 - \Phi^{\text{neg}}) \cdot h_j^{\text{neg}} \right)
\]

\[
\Delta b_2 = \eta_{\text{output}} \cdot \text{clamp}\left( (0 - \Phi^{\text{pos}}) + (1 - \Phi^{\text{neg}}),\; -1, 1 \right)
\]

The weight vector \(\mathbf{w}_2\) is renormalized to
\(\sum w_{2,j} = 1\) after each update to ensure stability. This
prevents runaway magnitude while preserving the relative distribution.

\paragraph{Why This Avoids Collapse}\label{why-this-avoids-collapse}

The hidden and output layers form a \textbf{competitive cooperative
system}:

\begin{itemize}
\tightlist
\item
  Units where \(h^{\text{neg}} > h^{\text{pos}}\) get a positive hidden
  update (h\_pos decreases further) and a positive w₂ update (their
  contribution to the output grows). These become \textbf{neg-selective
  detectors}.
\item
  Units where \(h^{\text{pos}} > h^{\text{neg}}\) get a positive hidden
  update (h\_pos increases further) but a negative w₂ update (their
  contribution shrinks). These units are \textbf{suppressed by the
  output layer}.
\end{itemize}

The output suppression dominates because w₂ can drop to zero,
eliminating a unit's contribution entirely. The system converges to a
state where high-w₂ units have \(h^{\text{neg}} \gg h^{\text{pos}}\),
yielding \(\Phi^{\text{pos}} \ll \Phi^{\text{neg}}\).

\begin{center}\rule{0.5\linewidth}{0.5pt}\end{center}

\subsection{3. Sequence Generation
Pipeline}\label{sequence-generation-pipeline}

\begin{enumerate}
\def\labelenumi{\arabic{enumi}.}
\tightlist
\item
  \textbf{Initialize}: Set the slow state \(s\) to a learned start
  vector (e.g., zero), and the fast state \(f\) to the prompt window.
\item
  \textbf{Evolve}: For each desired byte:

  \begin{itemize}
  \tightlist
  \item
    Run the Hamiltonian dynamics for \(K\) steps with the time‑dependent
    potential.
  \item
    Extract the last byte from the fast state and append it to the
    output.
  \item
    Shift the fast window: drop the first byte, append the new byte.
  \item
    Increment the time step \(t\) to increase the driving force.
  \end{itemize}
\item
  \textbf{Update slow state}: After a sentence boundary (detected by
  punctuation), update the slow state using a separate Hebbian rule (or
  simply average the fast state over the sentence).
\end{enumerate}

The system naturally produces variable‑length sequences because the
driving potential reduces energy for moving forward, but the damping and
the shifting operation ensure that the state doesn't become too
dispersed.

\begin{center}\rule{0.5\linewidth}{0.5pt}\end{center}

\subsection{4. Implementation Sketch
(C++)}\label{implementation-sketch-c}

\begin{Shaded}
\begin{Highlighting}[]
\CommentTok{// Hierarchical state}
\KeywordTok{struct}\NormalTok{ HierarchicalState }\OperatorTok{\{}
\NormalTok{    Vector s}\OperatorTok{;} \CommentTok{// slow (context)}
\NormalTok{    Vector f}\OperatorTok{;} \CommentTok{// fast (text window)}
\NormalTok{    Vector ps}\OperatorTok{,}\NormalTok{ pf}\OperatorTok{;} \CommentTok{// momenta}
\OperatorTok{\};}

\CommentTok{// Potential with flow term}
\DataTypeTok{double}\NormalTok{ potential}\OperatorTok{(}\AttributeTok{const}\NormalTok{ HierarchicalState}\OperatorTok{\&}\NormalTok{ x}\OperatorTok{,} \DataTypeTok{double}\NormalTok{ t}\OperatorTok{)} \OperatorTok{\{}
    \DataTypeTok{double}\NormalTok{ phi\_f }\OperatorTok{=}\NormalTok{ potential\_fast}\OperatorTok{(}\NormalTok{x}\OperatorTok{.}\NormalTok{f}\OperatorTok{,}\NormalTok{ x}\OperatorTok{.}\NormalTok{s}\OperatorTok{);}
    \DataTypeTok{double}\NormalTok{ phi\_s }\OperatorTok{=}\NormalTok{ potential\_slow}\OperatorTok{(}\NormalTok{x}\OperatorTok{.}\NormalTok{s}\OperatorTok{);}
    \DataTypeTok{double}\NormalTok{ phi\_flow }\OperatorTok{=} \OperatorTok{{-}}\NormalTok{alpha }\OperatorTok{*}\NormalTok{ dot}\OperatorTok{(}\NormalTok{gradient\_f}\OperatorTok{(}\NormalTok{phi\_f}\OperatorTok{,}\NormalTok{ x}\OperatorTok{.}\NormalTok{f}\OperatorTok{),}\NormalTok{ target\_flow}\OperatorTok{(}\NormalTok{x}\OperatorTok{.}\NormalTok{s}\OperatorTok{))} \OperatorTok{*}\NormalTok{ ramp}\OperatorTok{(}\NormalTok{t}\OperatorTok{);}
    \ControlFlowTok{return}\NormalTok{ phi\_f }\OperatorTok{+}\NormalTok{ phi\_s }\OperatorTok{+}\NormalTok{ phi\_flow}\OperatorTok{;}
\OperatorTok{\}}

\CommentTok{// Decoupled CHL: hidden and output layers updated separately}
\DataTypeTok{void}\NormalTok{ dchl\_update}\OperatorTok{(}\NormalTok{Potential}\OperatorTok{\&}\NormalTok{ pot}\OperatorTok{,} \AttributeTok{const}\NormalTok{ Vector}\OperatorTok{\&}\NormalTok{ c\_pos}\OperatorTok{,} \AttributeTok{const}\NormalTok{ Vector}\OperatorTok{\&}\NormalTok{ c\_neg}\OperatorTok{,}
                 \DataTypeTok{double}\NormalTok{ lr\_hidden}\OperatorTok{,} \DataTypeTok{double}\NormalTok{ lr\_output}\OperatorTok{)} \OperatorTok{\{}
\NormalTok{    Matrix h\_pos}\OperatorTok{,}\NormalTok{ h\_neg}\OperatorTok{;}
    \DataTypeTok{double}\NormalTok{ phi\_pos }\OperatorTok{=}\NormalTok{ pot}\OperatorTok{.}\NormalTok{forward}\OperatorTok{(}\NormalTok{c\_pos}\OperatorTok{,} \OperatorTok{\&}\NormalTok{h\_pos}\OperatorTok{);}
    \DataTypeTok{double}\NormalTok{ phi\_neg }\OperatorTok{=}\NormalTok{ pot}\OperatorTok{.}\NormalTok{forward}\OperatorTok{(}\NormalTok{c\_neg}\OperatorTok{,} \OperatorTok{\&}\NormalTok{h\_neg}\OperatorTok{);}

    \CommentTok{// Hidden: per{-}unit contrastive error (c\_pos only, self{-}amplifying)}
    \ControlFlowTok{for} \OperatorTok{(}\DataTypeTok{size\_t}\NormalTok{ j }\OperatorTok{=} \DecValTok{0}\OperatorTok{;}\NormalTok{ j }\OperatorTok{\textless{}}\NormalTok{ H}\OperatorTok{;} \OperatorTok{++}\NormalTok{j}\OperatorTok{)} \OperatorTok{\{}
        \DataTypeTok{double}\NormalTok{ error }\OperatorTok{=}\NormalTok{ h\_pos}\OperatorTok{[}\NormalTok{j}\OperatorTok{]} \OperatorTok{{-}}\NormalTok{ h\_neg}\OperatorTok{[}\NormalTok{j}\OperatorTok{];}
        \ControlFlowTok{for} \OperatorTok{(}\DataTypeTok{size\_t}\NormalTok{ i }\OperatorTok{=} \DecValTok{0}\OperatorTok{;}\NormalTok{ i }\OperatorTok{\textless{}}\NormalTok{ D}\OperatorTok{;} \OperatorTok{++}\NormalTok{i}\OperatorTok{)}
\NormalTok{            pot}\OperatorTok{.}\NormalTok{W1}\OperatorTok{(}\NormalTok{j}\OperatorTok{,}\NormalTok{ i}\OperatorTok{)} \OperatorTok{+=}\NormalTok{ lr\_hidden }\OperatorTok{*}\NormalTok{ error }\OperatorTok{*}\NormalTok{ c\_pos}\OperatorTok{[}\NormalTok{i}\OperatorTok{];}
\NormalTok{        pot}\OperatorTok{.}\NormalTok{b1}\OperatorTok{(}\NormalTok{j}\OperatorTok{)} \OperatorTok{+=}\NormalTok{ lr\_hidden }\OperatorTok{*}\NormalTok{ error}\OperatorTok{;}
    \OperatorTok{\}}

    \CommentTok{// Output: projection to targets, w2 renormalized to sum=1}
    \ControlFlowTok{for} \OperatorTok{(}\DataTypeTok{size\_t}\NormalTok{ j }\OperatorTok{=} \DecValTok{0}\OperatorTok{;}\NormalTok{ j }\OperatorTok{\textless{}}\NormalTok{ H}\OperatorTok{;} \OperatorTok{++}\NormalTok{j}\OperatorTok{)} \OperatorTok{\{}
\NormalTok{        pot}\OperatorTok{.}\NormalTok{w2}\OperatorTok{(}\NormalTok{j}\OperatorTok{)} \OperatorTok{+=}\NormalTok{ lr\_output }\OperatorTok{*} \OperatorTok{((}\DecValTok{0} \OperatorTok{{-}}\NormalTok{ phi\_pos}\OperatorTok{)} \OperatorTok{*}\NormalTok{ h\_pos}\OperatorTok{[}\NormalTok{j}\OperatorTok{]} \OperatorTok{+} \OperatorTok{(}\DecValTok{1} \OperatorTok{{-}}\NormalTok{ phi\_neg}\OperatorTok{)} \OperatorTok{*}\NormalTok{ h\_neg}\OperatorTok{[}\NormalTok{j}\OperatorTok{]);}
        \ControlFlowTok{if} \OperatorTok{(}\NormalTok{pot}\OperatorTok{.}\NormalTok{w2}\OperatorTok{(}\NormalTok{j}\OperatorTok{)} \OperatorTok{\textless{}} \DecValTok{0}\OperatorTok{)}\NormalTok{ pot}\OperatorTok{.}\NormalTok{w2}\OperatorTok{(}\NormalTok{j}\OperatorTok{)} \OperatorTok{=} \DecValTok{0}\OperatorTok{;}
    \OperatorTok{\}}
    \DataTypeTok{double}\NormalTok{ wsum }\OperatorTok{=}\NormalTok{ sum}\OperatorTok{(}\NormalTok{pot}\OperatorTok{.}\NormalTok{w2}\OperatorTok{);}
    \ControlFlowTok{if} \OperatorTok{(}\NormalTok{wsum }\OperatorTok{\textgreater{}} \DecValTok{0}\OperatorTok{)} \ControlFlowTok{for} \OperatorTok{(}\KeywordTok{auto}\OperatorTok{\&}\NormalTok{ w }\OperatorTok{:}\NormalTok{ pot}\OperatorTok{.}\NormalTok{w2}\OperatorTok{)}\NormalTok{ w }\OperatorTok{/=}\NormalTok{ wsum}\OperatorTok{;}

    \DataTypeTok{double}\NormalTok{ b2\_delta }\OperatorTok{=} \OperatorTok{(}\DecValTok{0} \OperatorTok{{-}}\NormalTok{ phi\_pos}\OperatorTok{)} \OperatorTok{+} \OperatorTok{(}\DecValTok{1} \OperatorTok{{-}}\NormalTok{ phi\_neg}\OperatorTok{);}
\NormalTok{    pot}\OperatorTok{.}\NormalTok{b2 }\OperatorTok{+=}\NormalTok{ lr\_output }\OperatorTok{*}\NormalTok{ clamp}\OperatorTok{(}\NormalTok{b2\_delta}\OperatorTok{,} \OperatorTok{{-}}\FloatTok{1.0}\OperatorTok{,} \FloatTok{1.0}\OperatorTok{);}
\OperatorTok{\}}

\CommentTok{// Generation}
\DataTypeTok{void}\NormalTok{ generate}\OperatorTok{(}\AttributeTok{const}\NormalTok{ string}\OperatorTok{\&}\NormalTok{ prompt}\OperatorTok{,} \DataTypeTok{int}\NormalTok{ max\_chars}\OperatorTok{)} \OperatorTok{\{}
\NormalTok{    HierarchicalState x}\OperatorTok{;}
\NormalTok{    x}\OperatorTok{.}\NormalTok{f }\OperatorTok{=}\NormalTok{ text\_to\_state}\OperatorTok{(}\NormalTok{prompt}\OperatorTok{);}
\NormalTok{    x}\OperatorTok{.}\NormalTok{s }\OperatorTok{=}\NormalTok{ Vector}\OperatorTok{(}\VariableTok{s\_dim}\OperatorTok{,} \FloatTok{0.0}\OperatorTok{);}
\NormalTok{    x}\OperatorTok{.}\NormalTok{pf }\OperatorTok{=}\NormalTok{ Vector}\OperatorTok{(}\NormalTok{f\_dim}\OperatorTok{,} \FloatTok{0.0}\OperatorTok{);}
\NormalTok{    x}\OperatorTok{.}\NormalTok{ps }\OperatorTok{=}\NormalTok{ Vector}\OperatorTok{(}\VariableTok{s\_dim}\OperatorTok{,} \FloatTok{0.0}\OperatorTok{);}
    \DataTypeTok{double}\NormalTok{ t }\OperatorTok{=} \FloatTok{0.0}\OperatorTok{;}
\NormalTok{    string output }\OperatorTok{=}\NormalTok{ prompt}\OperatorTok{;}
    \ControlFlowTok{for} \OperatorTok{(}\DataTypeTok{int}\NormalTok{ i }\OperatorTok{=} \DecValTok{0}\OperatorTok{;}\NormalTok{ i }\OperatorTok{\textless{}}\NormalTok{ max\_chars}\OperatorTok{;} \OperatorTok{++}\NormalTok{i}\OperatorTok{)} \OperatorTok{\{}
        \ControlFlowTok{for} \OperatorTok{(}\DataTypeTok{int}\NormalTok{ step }\OperatorTok{=} \DecValTok{0}\OperatorTok{;}\NormalTok{ step }\OperatorTok{\textless{}}\NormalTok{ K}\OperatorTok{;} \OperatorTok{++}\NormalTok{step}\OperatorTok{)} \OperatorTok{\{}
\NormalTok{            Vector grad }\OperatorTok{=}\NormalTok{ gradient}\OperatorTok{(}\NormalTok{potential}\OperatorTok{,}\NormalTok{ x}\OperatorTok{,}\NormalTok{ t}\OperatorTok{);}
\NormalTok{            x}\OperatorTok{.}\NormalTok{pf }\OperatorTok{{-}=} \FloatTok{0.5} \OperatorTok{*}\NormalTok{ dt }\OperatorTok{*}\NormalTok{ grad\_f}\OperatorTok{;}
\NormalTok{            x}\OperatorTok{.}\NormalTok{f }\OperatorTok{+=}\NormalTok{ dt }\OperatorTok{*}\NormalTok{ x}\OperatorTok{.}\NormalTok{pf}\OperatorTok{;}
\NormalTok{            x}\OperatorTok{.}\NormalTok{pf }\OperatorTok{{-}=} \FloatTok{0.5} \OperatorTok{*}\NormalTok{ dt }\OperatorTok{*}\NormalTok{ grad\_f}\OperatorTok{;}
\NormalTok{            x}\OperatorTok{.}\NormalTok{ps }\OperatorTok{{-}=} \FloatTok{0.5} \OperatorTok{*}\NormalTok{ dt }\OperatorTok{*}\NormalTok{ grad\_s}\OperatorTok{;}
\NormalTok{            x}\OperatorTok{.}\NormalTok{s }\OperatorTok{+=}\NormalTok{ dt }\OperatorTok{*}\NormalTok{ x}\OperatorTok{.}\NormalTok{ps}\OperatorTok{;}
\NormalTok{            x}\OperatorTok{.}\NormalTok{ps }\OperatorTok{{-}=} \FloatTok{0.5} \OperatorTok{*}\NormalTok{ dt }\OperatorTok{*}\NormalTok{ grad\_s}\OperatorTok{;}
        \OperatorTok{\}}
        \DataTypeTok{int}\NormalTok{ byte }\OperatorTok{=} \OperatorTok{(}\DataTypeTok{int}\OperatorTok{)(}\NormalTok{x}\OperatorTok{.}\NormalTok{f}\OperatorTok{[}\NormalTok{f\_dim}\OperatorTok{{-}}\DecValTok{1}\OperatorTok{]} \OperatorTok{*} \FloatTok{255.0} \OperatorTok{+} \FloatTok{0.5}\OperatorTok{);}
        \ControlFlowTok{if} \OperatorTok{(}\NormalTok{byte }\OperatorTok{==} \DecValTok{0}\OperatorTok{)} \ControlFlowTok{break}\OperatorTok{;}
\NormalTok{        output}\OperatorTok{.}\NormalTok{push\_back}\OperatorTok{((}\DataTypeTok{char}\OperatorTok{)}\NormalTok{byte}\OperatorTok{);}
\NormalTok{        shift\_window}\OperatorTok{(}\NormalTok{x}\OperatorTok{.}\NormalTok{f}\OperatorTok{,}\NormalTok{ byte}\OperatorTok{);}
\NormalTok{        t }\OperatorTok{+=}\NormalTok{ dt}\OperatorTok{;}
    \OperatorTok{\}}
\NormalTok{    cout }\OperatorTok{\textless{}\textless{}}\NormalTok{ output }\OperatorTok{\textless{}\textless{}}\NormalTok{ endl}\OperatorTok{;}
\OperatorTok{\}}
\end{Highlighting}
\end{Shaded}

\begin{center}\rule{0.5\linewidth}{0.5pt}\end{center}

\subsection{5. Why This Solves the
Problems}\label{why-this-solves-the-problems}

\begin{longtable}[]{@{}
  >{\raggedright\arraybackslash}p{(\linewidth - 4\tabcolsep) * \real{0.3333}}
  >{\raggedright\arraybackslash}p{(\linewidth - 4\tabcolsep) * \real{0.3333}}
  >{\raggedright\arraybackslash}p{(\linewidth - 4\tabcolsep) * \real{0.3333}}@{}}
\toprule\noalign{}
\begin{minipage}[b]{\linewidth}\raggedright
Problem
\end{minipage} & \begin{minipage}[b]{\linewidth}\raggedright
P‑CMA
\end{minipage} & \begin{minipage}[b]{\linewidth}\raggedright
HCMA Solution
\end{minipage} \\
\midrule\noalign{}
\endhead
\bottomrule\noalign{}
\endlastfoot
Hidden layer collapse & Scalar Φ global & \textbf{Decoupled CHL} with
per-unit contrastive errors + w₂ renormalization \\
Static potential trapping & Single particle, fixed Φ &
\textbf{Time‑dependent flow potential} with driving force \\
No sequence generation & Single minimum & \textbf{Hierarchical slow/fast
states} with forward bias \\
\end{longtable}

Crucially, HCMA retains all the advantages: O(1) memory, no
tokenization, no encoders/decoders, no loss functions, and no
hallucination.

\begin{center}\rule{0.5\linewidth}{0.5pt}\end{center}

\subsection{6. Conclusion}\label{conclusion}

We have introduced HCMA, a novel cognitive architecture that overcomes
the fundamental limitations of static potential‑based generation. By
decoupling the state into slow and fast components, adding a
time‑dependent driving potential, and employing decoupled contrastive
Hebbian learning with per-unit hidden errors and independent output
projection, HCMA enables coherent sequence generation while staying true
to the non‑Transformer, physics‑based philosophy. This work paves the
way for truly autonomous, memory‑efficient, and hallucination‑free AI
systems that can generate language naturally without ever resorting to
tokenization or autoregressive prediction.

\begin{center}\rule{0.5\linewidth}{0.5pt}\end{center}

\end{document}
